\section{ماژول فشرده‌سازی تحلیل مؤلفه‌های اصلی مستحکم (\lr{RPCA Compressor})}
\label{sec:rpca_compressor}

\subsection{معرفی و انگیزه}
این مکانیزم بر مبنای مبانی نظری پیشرفته در ادبیات توجه و فشرده‌سازی \cite{ref_rpcacompressor_2024} توسعه یافته است.

در فرآیند فشرده‌سازی و تجزیه ماتریس‌های وزن در مدل‌های زبانی بزرگ، یکی از چالش‌های بنیادین، وجود ناهنجاری‌ها و مقادیر پرت (\lr{Outliers}) در توزیع پارامترها و فعال‌سازی‌ها است. روش استاندارد تجزیه مقادیر منفرد (\lr{SVD}) به شدت در برابر مقادیر دورافتاده با دامنه بزرگ آسیب‌پذیر است؛ زیرا نرم فروبنیوس خطای بازسازی را با توان دوم محاسبه کرده و انحرافات مقطعی شدید، بخش عمده‌ای از انرژی رتبه‌های برتر \lr{SVD} را به خود اختصاص می‌دهند. این پدیده موجب تخریب زیرفضای معنایی اصلی لایه‌ها در نرخ‌های فشرده‌سازی بالا می‌گردد.

ماژول \lr{RPCA Compressor} با الهام از مبانی «تحلیل مؤلفه‌های اصلی مستحکم» (\lr{Robust Principal Component Analysis}) و برنامه‌ریزی تعقیب مؤلفه‌های اصلی (\lr{Principal Component Pursuit - PCP}) طراحی گردیده است. ایده بنیادین این الگوریتم، تفکیک صریح ماتریس وزن لایه خطی $\mathbf{W}$ به دو مؤلفه مجزا است:
\begin{enumerate}
    \item \textbf{مؤلفه کم‌رتبه ساختاری ($\mathbf{L}$):} حاوی همبستگی‌های اصلی و پایدار ماتریس وزن که قابلیت فشرده‌سازی با رتبه پایین را دارند.
    \item \textbf{مؤلفه تنک ناهنجاری‌ها و مقادیر پرت ($\mathbf{S}$):} حاوی انحرافات پراکنده و المان‌های کلیدی با دامنه بزرگ که به صورت یک ماتریس خلوت ایزوله می‌شوند.
\end{enumerate}
با جداسازی مؤلفه تنک $\mathbf{S}$، ماتریس پایه $\mathbf{L}$ به یک ساختار ریاضی همگن و خالص تبدیل گشته و اعمال تقریب کم‌رتبه بر روی آن، خطای بازسازی بسیار کمتری را نسبت به اعمال مستقیم \lr{SVD} بر روی ماتریس خام $\mathbf{W}$ به همراه می‌آورد.

\subsection{مبانی نظری و فرمول‌بندی ریاضی}
فرآیند تجزیه ماتریس وزن $\mathbf{W} \in \mathbb{R}^{d_{\text{out}} \times d_{\text{in}}}$ بر مبنای مسئله بهینه‌سازی محدب \lr{RPCA} به صورت زیر فرمول‌بندی می‌گردد:
\begin{equation}
\min_{\mathbf{L}, \mathbf{S}} \|\mathbf{L}\|_* + \lambda \|\mathbf{S}\|_1 \quad \text{subject to} \quad \mathbf{W} = \mathbf{L} + \mathbf{S}
\end{equation}
که در آن $\|\mathbf{L}\|_* = \sum_i \sigma_i(\mathbf{L})$ نرم هسته‌ای (تقریب محدب رتبه ماتریس) و $\|\mathbf{S}\|_1 = \sum_{i,j} |S_{ij}|$ نرم $\ell_1$ (تقریب محدب تَنَکی) است. پارامتر تنظیم‌کننده معمولاً به صورت نظری $\lambda = 1 / \sqrt{\max(d_{\text{out}}, d_{\text{in}})}$ تعیین می‌گردد.

در پیاده‌سازی بلادرنگ ماژول، جهت دستیابی به سرعت بالا در لایه‌های بزرگ شبکه، از الگوریتم بهینه‌سازی افزاینده لاگرانژ نامتعامد (\lr{Inexact ALM / PCP}) با عملگرهای تحلیلی آستانه‌گذاری در تکرارهای متناوب استفاده می‌شود:

\begin{enumerate}
    \item \textbf{گام تصویر کم‌رتبه با تجزیه مقادیر منفرد مقطوع (\lr{Singular Value Thresholding}):}
    در هر تکرار $t$، تجزیه مقادیر منفرد مؤلفه کم‌رتبه محاسبه شده و تصویر مقطوع با رتبه هدف $r$ بازسازی می‌گردد:
    \begin{equation}
    \mathbf{L}^{(t)} \approx \mathbf{U}_r \mathbf{\Sigma}_r \mathbf{V}_r^\top = \left(\mathbf{U}_{:, :r} \mathbf{\Sigma}_r^{1/2}\right) \left(\mathbf{\Sigma}_r^{1/2} \mathbf{V}_{:r, :}^\top\right)
    \end{equation}
    که در آن $r \le \min(d_{\text{out}}, d_{\text{in}} - 1)$ و سقف پارامتری $r \le \left\lfloor \frac{d_{\text{in}} d_{\text{out}}}{d_{\text{in}} + d_{\text{out}}} \right\rfloor - 1$ جهت تضمین کاهش قطعی حجم اعمال می‌شود.

    \item \textbf{محاسبه باقی‌مانده و آستانه‌گذاری نرم مؤلفه تنک (\lr{Soft-Thresholding}):}
    باقی‌مانده ماتریس نسبت به مؤلفه کم‌رتبه محاسبه شده و عملگر انقباض نرم (\lr{Shrinkage}) با آستانه $\mu$ بر آن اعمال می‌شود:
    \begin{equation}
    \mathbf{R}^{(t)} = \mathbf{W} - \mathbf{L}^{(t)}
    \end{equation}
    \begin{equation}
    \mathbf{S}^{(t)} = \mathcal{S}_{\mu}(\mathbf{R}^{(t)}) = \operatorname{sign}(\mathbf{R}^{(t)}) \odot \max(|\mathbf{R}^{(t)}| - \mu, 0)
    \end{equation}

    \item \textbf{به‌روزرسانی هدف کم‌رتبه:}
    مؤلفه تنک از ماتریس وزن اولیه تفریق شده و ماتریس کم‌رتبه برای گام بعدی بازآرایی می‌گردد:
    \begin{equation}
    \mathbf{L}^{(t+1)} = \mathbf{W} - \mathbf{S}^{(t)}
    \end{equation}
\end{enumerate}

پس از همگرایی تکرارها در سقف \lr{max\_iter}، ماتریس حاصله $\mathbf{L}$ از طریق تجزیه دومرحله‌ای به دو فاکتور وزن متوالی $\mathbf{A} \in \mathbb{R}^{r \times d_{\text{in}}}$ و $\mathbf{B} \in \mathbb{R}^{d_{\text{out}} \times r}$ تبدیل می‌شود:
\begin{equation}
\mathbf{A} = \mathbf{\Sigma}_r^{1/2} \mathbf{V}_r^\top, \quad \mathbf{B} = \mathbf{U}_r \mathbf{\Sigma}_r^{1/2} \implies \mathbf{y} = \mathbf{x} \mathbf{A}^\top \mathbf{B}^\top + \mathbf{b}
\end{equation}
در صورت فعال بودن شاخه انطباق سریع، ماتریس تصحیح رتبه-پایین کمکی (\lr{LoRA}) نیز به ساختار لایه فشرده متصل می‌گردد.

\paragraph{تحلیل طیفی و نامساوی‌های پایداری خطا}
برای بررسی میزان دقت تقریب ماتریس وزن یا نگاشت توجه در ماژول \lr{rpca\_compressor}، تفاضل ماتریس فاکتورگیری‌شده از نگاشت متراکم اصلی بر اساس نامساوی هافدینگ-فون‌نیومان ارزیابی می‌گردد. فرض کنید $\mathbf{W} \in \mathbb{R}^{d_{\text{in}} \times d_{\text{out}}}$ ماتریس وزن هدف باشد که توسط تقریب کم‌مرتبه یا پراکنده با رتبه موثر $r$ جایگزین شده است. کران بالای خطای بازسازی بر مبنای نرم فروبنیوس به صورت زیر تصریح می‌شود:
\begin{equation}
\label{eq:spectral_bound_rpca_compressor}
\| \mathbf{W} - \mathbf{W}_{r} \|_F = \sqrt{\sum_{i=r+1}^{\min(d_{\text{in}}, d_{\text{out}})} \sigma_i^2(\mathbf{W})} \le \epsilon \| \mathbf{W} \|_F
\end{equation}
که در آن $\sigma_i(\mathbf{W})$ بیانگر $i$-امین مقدار منفرد مرتب‌شده ماتریس به صورت نزولی است. نسبت حفظ انرژی طیفی $\eta_r = \frac{\sum_{i=1}^r \sigma_i^2}{\sum_{i=1}^{d} \sigma_i^2}$ تضمین می‌کند که با انتخاب آستانه $\tau \ge 0.90$، بیش از نود درصد از پراکندگی واریانس داده‌ها حفظ شده و افت سیگنال در انتشار گرادیان رو به عقب به حداقل ممکن کاهش می‌یابد.

\paragraph{تحلیل مجانبی و پیچیدگی اشتراک حافظه سخت‌افزاری}
از منظر پیچیدگی محاسباتی و انتقال داده‌ها در سلسومراتب حافظه (حافظه سراسری \lr{HBM} به حافظه محلی \lr{SRAM})، عملیات ضرب ماتریسی در این مکانیزم نیازمند تخصیص بافر موقت با ابعاد فاکتورگیری‌شده است. برای پردازش دنباله‌ای به طول $N$ با بعد مخفی $D$، بار محاسباتی هسته از فرم درجه دو به رابطه زیر تقلیل می‌یابد:
\begin{equation}
\label{eq:complexity_rpca_compressor}
\mathcal{C}_{\text{rpca\_compressor}} = \mathcal{O}\left( B \cdot N \cdot D \cdot r + B \cdot H \cdot N \cdot d_k \right) \ll \mathcal{O}\left( B \cdot H \cdot N^2 \cdot d_k \right)
\end{equation}
این بهبود مرتبه پیچیدگی، با کاهش چشمگیر حجم تبادل تانسورها در گذرگاه‌های ارتباطی کارت گرافیک، مانع از اشباع پهنای باند حافظه (پدیده \lr{Memory-Bound}) در فرآیند استنتاج طولانی‌مدت می‌گردد.

\subsection{معماری و پیاده‌سازی}
پیاده‌سازی این الگوریتم در مسیر \lr{Codes/framework/compression\_engine/rpca\_compressor.py} قرار داشته و با نام‌های \lr{rpca} و \lr{rpca\_compressor} در رجیستری فشرده‌سازی فراخوانی می‌شود. از ویژگی‌های بارز معماری و مهندسی این ماژول می‌توان به موارد ذیل اشاره کرد:
\begin{itemize}
    \item \textbf{محاسبات پایدار با دقت \lr{FP32}:} تمامی مراحل چرخه‌های تکرار \lr{PCP} و تجزیه‌های \lr{torch.linalg.svd} در فرمت ممیز شناور ۳۲ بیتی محاسبه شده و مقادیر منفرد با کف $10^{-12}$ جهت ممانعت از ناپایداری عددی محدود می‌شوند.
    \item \textbf{سازگاری کامل با لایه‌های خطی و کانوولوشن ۱ بعدی (\lr{Conv1D}):} تشخیص خودکار لایه‌های کانوولوشن ۱ بعدی (مانند ساختار \lr{GPT-2}) از طریق وارسی نام کلاس و ویژگی \lr{nf}، ترانهش صحیح ماتریس‌ها و نگاشت ابعاد را تضمین می‌کند.
    \item \textbf{مکانیزم تنزیل اضطراری دومرحله‌ای (\lr{Robust Fallback Hierarchy}):} در صورت بروز خطا در چرخه‌های \lr{ALM}، الگوریتم ابتدا به \lr{SVD} مستقیم روی ماتریس خام رجوع کرده و در صورت عدم دستیابی به شرط پایداری، لایه اصلی را بدون تخریب حفظ می‌نماید.
    \item \textbf{پایش تفکیکی خطای فروبنیوس:} خطای بازسازی نسبی فروبنیوس $\epsilon_{\text{recon}} = \frac{\|\mathbf{W} - \mathbf{B} \mathbf{A}\|_F}{\|\mathbf{W}\|_F}$ برای تک‌تک لایه‌ها ثبت و پایش می‌شود.
\end{itemize}

\subsection{نتایج تجربی}
ارزیابی عملکرد ماژول فشرده‌ساز \lr{RPCA Compressor} بر روی مدل‌های زبانی استاندارد در پایگاه داده بنچ‌مارک در جدول \ref{tab:benchmark_rpca_compressor} گزارش شده است:

\begin{table}[H]
\centering
\small
\caption{نتایج بنچ‌مارک فشرده‌ساز مبتنی بر تحلیل مؤلفه‌های اصلی مستحکم (\lr{RPCA Compressor})}
\label{tab:benchmark_rpca_compressor}
\begin{tabular}{lccccc}
\toprule
\textbf{مدل پایه \lr{(Base Model)}} & \textbf{دقت \lr{(Accuracy)}} & \textbf{شاخص \lr{F1}} & \textbf{حافظه \lr{(Peak VRAM)}} & \textbf{توان \lr{(Throughput)}} & \textbf{وضعیت} \\
\midrule
\lr{DistilBERT} & \lr{90.37\%} & \lr{0.9067} & \lr{311.19 MB} & \lr{138,163.61 tok/s} & \lr{PASS} \\
\lr{RoBERTa} & \lr{72.02\%} & \lr{0.7724} & \lr{301.72 MB} & \lr{92,321.44 tok/s} & \lr{PASS} \\
\lr{GPT-2} & \lr{60.21\%} & \lr{0.6825} & \lr{637.67 MB} & \lr{65,063.68 tok/s} & \lr{PASS} \\
\lr{Qwen} & -- & -- & -- & -- & \lr{SKIP} \\
\bottomrule
\end{tabular}
\end{table}

بررسی داده‌های عددی بنچ‌مارک نشان‌دهنده استحکام چشمگیر این ماژول در تمامی مدل‌های ارزیابی‌شده است:
\begin{itemize}
    \item در مدل \lr{DistilBERT}، دقت مدل از مقدار پایه \lr{90.48\%} به \lr{90.37\%} رسیده است که تنها بیانگر افت ناچیز \lr{0.11\%} بوده و هم‌زمان توان عملیاتی قابل توجه \lr{138,163.61 tok/s} را با وضعیت \lr{PASS} ثبت نموده است.
    \item در مدل \lr{RoBERTa}، دقت \lr{72.02\%} با مصرف حافظه بهینه \lr{301.72 MB} و توان عملیاتی \lr{92,321.44 tok/s} به دست آمده و افت دقت در محدوده مجاز کنترل شده است.
    \item در مدل زبانی اتورگرسیو \lr{GPT-2}، دقت \lr{60.21\%} و شاخص \lr{F1} برابر با \lr{0.6825} با توان گذردهی بیش از ۶۵ هزار توکن بر ثانیه احراز شده است.
    \item آزمایش روی مدل \lr{Qwen} به دلیل محدودیت منابع حافظه اختصاص‌یافته در اسکریپت یکپارچه در این نسخه در وضعیت عدم اجرا (\lr{SKIP}) قرار گرفته است.
\end{itemize}

\subsection{محدودیت‌ها و تحلیل شکست}
با وجود نتایج پایدار، تحلیل رفتار محاسباتی ماژول \lr{RPCA Compressor} مشخص‌کننده محدودیت‌های فنی زیر است:
\begin{itemize}
    \item \textbf{سربار زمانی چرخه‌های \lr{SVD} در فاز آفلاین:} اجرای چندباره عملگر \lr{SVD} در هر گام از تکرارهای \lr{ALM} برای لایه‌های حجیم (نظیر لایه‌های پیش‌خوران \lr{MLP}) زمان پیش‌پردازش کالیبراسیون را به صورت خطی با تعداد تکرارها افزایش می‌دهد.
    \item \textbf{حساسیت هایپرپارامتر آستانه‌گذاری $\mu$:} چنانچه پارامتر $\mu$ بیش از حد کوچک انتخاب شود، مؤلفه تنک $\mathbf{S}$ المان‌های غیرپرت را نیز جذب کرده و منجر به تضعیف اطلاعات مفید در مؤلفه کم‌رتبه $\mathbf{L}$ می‌گردد؛ در مقابل، مقادیر بسیار بزرگ $\mu$ کارایی جداسازی ناهنجاری‌ها را مختل می‌سازد.
\end{itemize}
